\documentclass[../main]{subfiles}
\begin{document}

\chapter{Observación de Modelos y Prototipos Hidrodinámicamente Estables}

\section{Introducción a la Hidrodinámica}

La hidrodinámica es la rama de la física que estudia el movimiento de los fluidos y las fuerzas que actúan sobre ellos. Es fundamental en el diseño de modelos y prototipos que deben ser estables cuando interactúan con fluidos en movimiento.

\section{Principios Básicos de la Hidrodinámica}

\subsection{Ecuaciones Fundamentales}

Existen varias ecuaciones fundamentales en la hidrodinámica, pero algunas de las más importantes para entender la estabilidad hidrodinámica son las siguientes:

\info{Ecuación de Continuidad}{
La ecuación de continuidad establece que, para un fluido incompresible, la masa que entra en un volumen debe ser igual a la masa que sale. Matemáticamente, se expresa como:
\begin{equation}
\nabla \cdot \mathbf{v} = 0
\end{equation}

Siendo:
\begin{itemize}
    \item $\nabla \cdot \mathbf{v}$: Divergencia del vector velocidad $\mathbf{v}$\footnote{La divergencia de un campo vectorial es una medida de la densidad neta de flujo saliendo de un punto.}.
\end{itemize}
}

\info{Ecuación de Bernoulli}{
La ecuación de Bernoulli es una de las más utilizadas en hidrodinámica y relaciona la presión, la velocidad y la altura en un flujo de fluido ideal. Se expresa como:
\begin{equation}
P + \frac{1}{2}\rho v^2 + \rho gh = \text{constante}
\end{equation}

Siendo:
\begin{itemize}
    \item $P$: Presión del fluido [Pa]
    \item $\rho$: Densidad del fluido [kg/m$^3$]
    \item $v$: Velocidad del fluido [m/s]
    \item $g$: Aceleración debida a la gravedad [m/s$^2$]
    \item $h$: Altura del fluido [m]
\end{itemize}
}

\info{Ecuación de Navier-Stokes}{
Las ecuaciones de Navier-Stokes describen el movimiento de los fluidos viscosos y son fundamentales para estudiar la estabilidad hidrodinámica. En su forma más simple, una de las ecuaciones se puede escribir como:
\begin{equation}
\rho \left( \frac{\partial \mathbf{v}}{\partial t} + \mathbf{v} \cdot \nabla \mathbf{v} \right) = -\nabla P + \mu \nabla^2 \mathbf{v} + \mathbf{f}
\end{equation}

Siendo:
\begin{itemize}
    \item $\rho$: Densidad del fluido [kg/m$^3$]
    \item $\mathbf{v}$: Vector velocidad del fluido [m/s]
    \item $P$: Presión [Pa]
    \item $\mu$: Viscosidad dinámica [Pa·s]
    \item $\mathbf{f}$: Fuerza externa por unidad de volumen [N/m$^3$]
\end{itemize}
}

\subsection{Estabilidad Hidrodinámica}

La estabilidad hidrodinámica se refiere a la capacidad de un cuerpo sumergido o flotante de regresar a su posición de equilibrio después de una perturbación. Esto es crucial en el diseño de barcos, submarinos y otros modelos que interactúan con fluidos.

\section{Ejemplo de Cálculo: Aplicación de la Ecuación de Bernoulli}

\ejemplo{Ejemplo}{Calcular la velocidad del agua que fluye a través de una tubería con una diferencia de presión entre dos puntos de 10,000 Pa. La densidad del agua es 1000 kg/m$^3$ y se considera que la altura se mantiene constante.

Datos:
\begin{itemize}
    \item $P_1 = 10,000$ Pa
    \item $P_2 = 0$ Pa (Presión de referencia)
    \item $\rho = 1000$ kg/m$^3$
    \item $h_1 = h_2$
\end{itemize}

Aplicamos la ecuación de Bernoulli:
\begin{equation}
P_1 + \frac{1}{2} \rho v_1^2 = P_2 + \frac{1}{2} \rho v_2^2
\end{equation}

Asumiendo que la velocidad inicial $v_1$ es pequeña y despreciable:
\begin{equation}
10,000 + 0 \approx 0 + \frac{1}{2} \cdot 1000 \cdot v_2^2
\end{equation}

Resolviendo para $v_2$:
\begin{equation}
v_2 = \sqrt{\frac{2 \cdot 10,000}{1000}} = \sqrt{20} \approx 4.47 \text{ m/s}
\end{equation}
}

\section{Cuestionario}

\actividad{Responder el siguiente cuestionario}
\begin{enumerate}
    \item ¿Qué es la ecuación de continuidad y cuál es su importancia en hidrodinámica?
    \item ¿Cómo se aplica la ecuación de Bernoulli en el diseño de prototipos hidrodinámicos?
    \item ¿Qué representan las ecuaciones de Navier-Stokes y por qué son fundamentales en hidrodinámica?
    \item ¿Cómo se define la estabilidad hidrodinámica?
    \item ¿Cuál es la relación entre la presión y la velocidad en un flujo de fluido según Bernoulli?
    \item Explica cómo la viscosidad afecta el movimiento de los fluidos.
    \item ¿Qué factores influyen en la estabilidad de un barco en el agua?
    \item Describe un ejemplo real donde se aplique la hidrodinámica en ingeniería.
    \item ¿Cómo se mide la densidad de un fluido y por qué es importante?
    \item ¿Qué papel juega la gravedad en la hidrodinámica?
\end{enumerate}

\section{Problemas}

\actividad{Resolver los siguientes problemas}
\begin{enumerate}
    \item Calcular la presión en un punto de una tubería donde el agua fluye a 3 m/s, si en otro punto de la tubería la presión es 5000 Pa y la velocidad es 1 m/s. La densidad del agua es 1000 kg/m$^3$.
    \item Determinar la velocidad de un fluido que sale de un orificio en un tanque a 2 metros de altura sobre el nivel del agua. Considerar la densidad del fluido como 1000 kg/m$^3$.
    \item Encontrar la altura a la que un fluido debe ser almacenado para que salga a una velocidad de 5 m/s por un orificio. La densidad del fluido es 1000 kg/m$^3$.
    \item Calcular la fuerza externa necesaria para mantener un flujo constante en una tubería donde la velocidad del fluido es 2 m/s y la densidad es 1000 kg/m$^3$. La viscosidad es 0.001 Pa·s.
    \item Determinar la velocidad de un fluido en una tubería si la presión en un punto es 2000 Pa y en otro punto es 0 Pa, con una densidad de 1000 kg/m$^3$.
    \item Calcular la diferencia de presión en una tubería de agua con una densidad de 1000 kg/m$^3$ y una velocidad de flujo de 4 m/s.
    \item Encontrar la presión en un punto de una tubería si la velocidad del fluido es 3 m/s y la densidad del fluido es 1000 kg/m$^3$, considerando que la presión en otro punto es 10000 Pa y la velocidad es 2 m/s.
    \item Determinar la velocidad de un fluido en una tubería si la presión inicial es 15000 Pa y la presión final es 5000 Pa, con una densidad de 1000 kg/m$^3$.
    \item Calcular la altura a la que un fluido debe ser almacenado para que salga a una velocidad de 6 m/s por un orificio, con una densidad de 1000 kg/m$^3$.
    \item Encontrar la velocidad de un fluido en una tubería donde la presión inicial es 8000 Pa y la presión final es 2000 Pa, con una densidad de 1000 kg/m$^3$.
\end{enumerate}

\end{document}
